\subsection*{Декларативное описание модели}

\begin{enumerate}
    \item[\textbf{а)}] Пользователь задаёт векторы $A$ и $B$ (режим генерации или ручной ввод через параметры/интерактивно).
    
    \item[\textbf{б)}] Если векторы имеют разный размер или числа выходят за диапазон разрядности $p=4$ (допустимы значения 0--15), модель останавливает работу и сообщает об ошибке.
    
    \item[\textbf{в)}] Формируется список пар $(A[i], B[i])$ для попарного умножения; при выводе пары могут отображаться в десятичном и при необходимости в двоичном виде (группами по 4 символа).
    
    \item[\textbf{г)}] Конвейер состоит из $n=4$ стадий. На каждой стадии хранятся: индекс обрабатываемой пары, множимое $A[i]$, множитель $B[i]$ и частичная сумма (8 бит). В начале такта на первой стадии частичная сумма равна нулю; на последующих стадиях --- результат предыдущей стадии.
    
    \item[\textbf{д)}] Для каждого такта модель определяет состояние конвейера: на стадии $s$ в такт $t$ обрабатывается пара с индексом $t-s$ (если он в диапазоне $[0, m)$). В первый такт в первую стадию поступает пара $[0]$; на каждом следующем такте в первую стадию поступает следующая пара из списка, пока пары не закончатся.
    
    \item[\textbf{е)}] Данные, вышедшие из последней (четвёртой) стадии, дают готовое произведение --- значение частичной суммы после четырёх шагов алгоритма есть результат умножения пары.
    
    \item[\textbf{ж)}] На каждой стадии выполняется один шаг алгоритма с младших разрядов со сдвигом частичной суммы вправо: если текущий разряд множителя (с номером, соответствующим стадии) равен 1, в старшую половину 8-битного регистра прибавляется множимое ($a \ll 4$); затем частичная сумма сдвигается вправо на 1 разряд. Результат передаётся на следующую стадию (в следующий такт).
    
    \item[\textbf{з)}] Если обработаны все такты (последний такт $t = m + n - 1$, когда из конвейера вышла последняя пара), модель завершает работу. Иначе выполняется переход к следующему такту (шаги д--з).
\end{enumerate}

\subsection*{Описание входов и выходов программы}

\textbf{Вход:} векторы $A$ и $B$ из 4-разрядных чисел (0--15); параметры конвейера: число стадий $n=4$, число пар $m$, время такта $t_i$ (усл. ед.). Ввод --- генерация (с опцией seed) или ручной (через аргументы или с клавиатуры).

\textbf{Выход:} протокол работы конвейера (пары до первой стадии; по тактам --- для каждой стадии индекс пары, частичная сумма в двоичном виде, время; после последней стадии --- элементы результирующего вектора $C$ с временем получения); краткая сводка (векторы $A$, $B$, $C$).

\subsection*{Исходные данные}

\begin{itemize}
    \item Разрядность $p$ попарно умножаемых чисел --- 4.
    \item Разрядность результата умножения --- $2p = 8$.
    \item Количество стадий конвейера $n$ --- 4 (равно $p$).
    \item Количество пар $m$ задаётся пользователем или при генерации (длина векторов).
    \item Время одного такта $t_i$ задаётся параметром (по умолчанию 1 усл. ед.).
\end{itemize}
